%% Generated by Sphinx.
\def\sphinxdocclass{jupyterBook}
\documentclass[letterpaper,10pt,english]{jupyterBook}
\ifdefined\pdfpxdimen
   \let\sphinxpxdimen\pdfpxdimen\else\newdimen\sphinxpxdimen
\fi \sphinxpxdimen=.75bp\relax
\ifdefined\pdfimageresolution
    \pdfimageresolution= \numexpr \dimexpr1in\relax/\sphinxpxdimen\relax
\fi
%% let collapsible pdf bookmarks panel have high depth per default
\PassOptionsToPackage{bookmarksdepth=5}{hyperref}
%% turn off hyperref patch of \index as sphinx.xdy xindy module takes care of
%% suitable \hyperpage mark-up, working around hyperref-xindy incompatibility
\PassOptionsToPackage{hyperindex=false}{hyperref}
%% memoir class requires extra handling
\makeatletter\@ifclassloaded{memoir}
{\ifdefined\memhyperindexfalse\memhyperindexfalse\fi}{}\makeatother

\PassOptionsToPackage{booktabs}{sphinx}
\PassOptionsToPackage{colorrows}{sphinx}

\PassOptionsToPackage{warn}{textcomp}

\catcode`^^^^00a0\active\protected\def^^^^00a0{\leavevmode\nobreak\ }
\usepackage{cmap}
\usepackage{fontspec}
\defaultfontfeatures[\rmfamily,\sffamily,\ttfamily]{}
\usepackage{amsmath,amssymb,amstext}
\usepackage{polyglossia}
\setmainlanguage{english}



\setmainfont{FreeSerif}[
  Extension      = .otf,
  UprightFont    = *,
  ItalicFont     = *Italic,
  BoldFont       = *Bold,
  BoldItalicFont = *BoldItalic
]
\setsansfont{FreeSans}[
  Extension      = .otf,
  UprightFont    = *,
  ItalicFont     = *Oblique,
  BoldFont       = *Bold,
  BoldItalicFont = *BoldOblique,
]
\setmonofont{FreeMono}[
  Extension      = .otf,
  UprightFont    = *,
  ItalicFont     = *Oblique,
  BoldFont       = *Bold,
  BoldItalicFont = *BoldOblique,
]



\usepackage[Bjarne]{fncychap}
\usepackage[,numfigreset=1,mathnumfig]{sphinx}

\fvset{fontsize=\small}
\usepackage{geometry}


% Include hyperref last.
\usepackage{hyperref}
% Fix anchor placement for figures with captions.
\usepackage{hypcap}% it must be loaded after hyperref.
% Set up styles of URL: it should be placed after hyperref.
\urlstyle{same}

\addto\captionsenglish{\renewcommand{\contentsname}{Bericht zur Tagung}}

\usepackage{sphinxmessages}



        % Start of preamble defined in sphinx-jupyterbook-latex %
         \usepackage[Latin,Greek]{ucharclasses}
        \usepackage{unicode-math}
        % fixing title of the toc
        \addto\captionsenglish{\renewcommand{\contentsname}{Contents}}
        \hypersetup{
            pdfencoding=auto,
            psdextra
        }
        % End of preamble defined in sphinx-jupyterbook-latex %
        

\title{Nachbericht DHd}
\date{Mar 17, 2026}
\release{}
\author{<a href="https://github.com/dh\sphinxhyphen{}erfurt">Nina Brolich, Digital Humanities Erfurt</a>}
\newcommand{\sphinxlogo}{\vbox{}}
\renewcommand{\releasename}{}
\makeindex
\begin{document}

\pagestyle{empty}
\sphinxmaketitle
\pagestyle{plain}
\sphinxtableofcontents
\pagestyle{normal}
\phantomsection\label{\detokenize{home::doc}}


\sphinxAtStartPar
Auch für die diesjährige DHd\sphinxhyphen{}Konferenz in Wien schrieb der DHd\sphinxhyphen{}Verband, in Zusammenarbeit mit \sphinxstyleemphasis{CLARIAH\sphinxhyphen{}AT} und den Konsortien \sphinxstyleemphasis{NFDI4Culture} und \sphinxstyleemphasis{NFDI4Memory}, wieder \sphinxhref{https://dhd-blog.org/?p=23004}{\sphinxstyleemphasis{Early Career Reisestipendien}} aus. Teil des Stipendienprogrammes ist ein Konferenzbegleitungsbeitrag, in dem die Stipendiat*innen von ihren Erfahrungen auf der Tagung berichten. Ich freue mich sehr, dass auch ich eines dieser Stipendien zugesprochen bekommen habe und demnach hier meinen Konferenzbegleitungsbeitrag in Form eines JupyterBooks präsentieren darf. Herzlich bedanken möchte ich mich insbesondere bei \sphinxstyleemphasis{NFDI4Memory}, die mein Stipendium finanziert haben.

\begin{DUlineblock}{0em}
\item[] \sphinxstylestrong{\Large Kurze Vorstellung}
\end{DUlineblock}

\sphinxAtStartPar
Ich heiße Nina Brolich und bin seit April 2025 wissenschaftliche Mitarbeiterin an der Professur für Digital Humanities, sowohl an der Fachhochschule als auch an der \sphinxhref{https://www.uni-erfurt.de/philosophische-fakultaet/seminare-professuren/historisches-seminar/professuren/digital-humanities-hybride-bildungs-und-kommunikationsraeume}{Universität Erfurt}. Besonders interessiere ich mich neben Digitalen Editionen für die Anwendung von Machine Learning in den Geisteswissenschaften und computerlinguistische Fragestellungen. Als Early\sphinxhyphen{}Career\sphinxhyphen{}DH\sphinxhyphen{}Wissenschaftlerin stehe ich nun vor der Aufgabe, mich innerhalb der Disziplin wissenschaftlich zu verorten und mir ein eigenständiges
Forschungsfeld zu erschließen. Die Tagung bietet für mich einen idealen Rahmen, um
meine Interessen im Austausch mit der DH\sphinxhyphen{}Community weiter zu konkretisieren und auch zu
erwägen, in welcher Fachgruppe des DHd\sphinxhyphen{}Verbands ich mich einbringen kann.

\begin{DUlineblock}{0em}
\item[] \sphinxstylestrong{\large Warum das Ganze nun als JupyterBook und nicht etwa als Blogbeitrag o.ä.?}
\end{DUlineblock}

\sphinxAtStartPar
Diese Konferenzbegleitung wird bewusst als JupyterBook umgesetzt. Ziel ist es, Inhalte strukturiert und nachvollziehbar aufzubereiten. Ein JupyterBook ermöglicht die klare Gliederung komplexer Themen in logisch aufeinander aufbauende Kapitel und Unterabschnitte und verbindet dabei Text, Abbildungen, Daten und – wo sinnvoll – ausführbaren Code in einer gemeinsamen Umgebung. Zugleich ist die Anwendung vergleichsweise niedrigschwellig und gut dokumentiert. Als etabliertes Format in der wissenschaftlichen Praxis, vor allem im Bereich von Open Educational Resources (OER), steht es für Transparenz, Nachnutzbarkeit und Offenheit. Gerade deshalb bietet es sich an, auch ein Konferenzformat in dieser Form zu erproben und damit wissenschaftliche Kommunikation nicht nur inhaltlich, sondern auch strukturell offen zu gestalten.

\begin{DUlineblock}{0em}
\item[] \sphinxstylestrong{\large Aufbau}
\end{DUlineblock}

\sphinxAtStartPar
Das JupyterBook ist thematisch in drei Hauptteile gegliedert.
\begin{itemize}
\item {} 
\sphinxAtStartPar
Der erste Teil, „Bericht zur Tagung“, bündelt Beiträge zu einzelnen thematischen Schwerpunkten wie Machine Learning und digitalen Editionen, gibt aber auch Einblick in andere Vorträge, die ich besucht habe.

\item {} 
\sphinxAtStartPar
Der zweite Teil widmet sich dem Mentoring\sphinxhyphen{}Programm und kombiniert eine allgemeine Programmbeschreibung mit einer persönlichen Reflexion meiner Teilnahme am Programm.

\item {} 
\sphinxAtStartPar
Der dritte Teil richtet den Blick auf die Arbeitsgruppen des DHd\sphinxhyphen{}Verbandes.

\end{itemize}

\sphinxAtStartPar
Stand: März 2026

\sphinxstepscope


\part{Bericht zur Tagung}

\sphinxstepscope


\chapter{Machine Learning}
\label{\detokenize{tagung/machine-learning:machine-learning}}\label{\detokenize{tagung/machine-learning::doc}}
\sphinxAtStartPar
Machine\sphinxhyphen{}Learning\sphinxhyphen{}Verfahren durchziehen zunehmend auch die geisteswissenschaftliche Forschung – sei es in der Bildanalyse, der Textklassifikation oder der Sentiment\sphinxhyphen{}Analyse. Gerade im Kontext des anhaltenden LLM\sphinxhyphen{}Booms bot die DHd zahlreiche Beiträge, die konkrete Anwendungen erproben oder die Verwendung methodisch reflektieren. Für dieses Kapitel habe ich zwei Vorträge ausgewählt, die sich der automatisierten Extraktion und Erschließung historischer Quellen widmen – ein Thema, das meiner eigenen Forschungsarbeit besonders nahesteht.


\section{Doctoral Consortium: Von der Handschrift zum Datensatz: Computergestützte Erschließung und Aufbereitung historischer Wetterdaten (Constantin Lehenmeier)}
\label{\detokenize{tagung/machine-learning:doctoral-consortium-von-der-handschrift-zum-datensatz-computergestutzte-erschlieszung-und-aufbereitung-historischer-wetterdaten-constantin-lehenmeier}}
\sphinxAtStartPar
Constantin Lehenmeier stellte sein Dissertationsprojekt zur computergestützten Erschließung historischer Wetterdaten vor, wobei die Modellierung von Tabellen als spezifische Datenstrukturen im Mittelpunkt stand – ein oft unterschätztes Problem, das er mit einer Kombination aus OCR\sphinxhyphen{}D, einem eigens trainierten YOLO\sphinxhyphen{}Modell und NER\sphinxhyphen{}Verfahren (HmBERT, GND\sphinxhyphen{}Verlinkung) angeht. Besonders interessant war die kritische Reflexion quantitativer Metriken: Sie validieren Modelle, nicht historische Korrektheit. Hervorzuheben ist außerdem der bewusste Einsatz von Open\sphinxhyphen{}Source\sphinxhyphen{}Software und \sphinxhyphen{}Modellen als konsequente Entscheidung für Verfügbarkeit, Nachnutzbarkeit und langfristige Nachhaltigkeit der erzeugten Ressourcen.


\section{Automatisierte Datenextraktion im Rahmen des FWF\sphinxhyphen{}Projekts ‚Digitale Erschließung des Schematismus‘ (Bernhard Ortbauer)}
\label{\detokenize{tagung/machine-learning:automatisierte-datenextraktion-im-rahmen-des-fwf-projekts-digitale-erschlieszung-des-schematismus-bernhard-ortbauer}}
\sphinxAtStartPar
Bernhard Ortbauer stellte die automatisierte Datenextraktion im Rahmen des FWF\sphinxhyphen{}Projekts \sphinxstyleemphasis{Digitale Erschließung des Schematismus} vor, bei der Layout\sphinxhyphen{}Erkennung (YOLO, Transformer) und NER kombiniert werden, um die hierarchischen Informationsstrukturen des Drucks in einen Knowledge Graph zu überführen. Dabei setzt er sich mit Problemstellungen wie visueller Ähnlichkeit auseinander: Eindeutige Lösungen gibt es nicht – stattdessen müssen Regelmäßigkeiten systematisch ausgenutzt und semantische Informationen gezielt, aber sparsam einbezogen werden.

\sphinxstepscope


\chapter{Digitale Editionen}
\label{\detokenize{tagung/digitale-editionen:digitale-editionen}}\label{\detokenize{tagung/digitale-editionen::doc}}
\sphinxAtStartPar
Digitale Editionen gehören zu den etabliertesten Feldern der Digital Humanities – und doch bleibt die Frage, wie Texte erschlossen, annotiert und publiziert werden sollten, stets neu verhandelbar. Für dieses Kapitel habe ich zwei Beiträge ausgewählt, die das Feld der digitalen Editionen aus sehr unterschiedlichen Blickwinkeln beleuchten: einerseits ein hands\sphinxhyphen{}on Workshop zu einem etablierten Publikationswerkzeug, andererseits ein konzeptueller Beitrag, der die medialen Grenzen der Edition selbst in den Blick nimmt.


\section{Workshop: TEI Publisher reloaded: Digitale Editionen mit System – modular, nachhaltig, community\sphinxhyphen{}orientiert (Richard Diebel, Nils Berns, Andreas Christ, Wolfgang Meier, Magdalena Turska, Lars Windauer)}
\label{\detokenize{tagung/digitale-editionen:workshop-tei-publisher-reloaded-digitale-editionen-mit-system-modular-nachhaltig-community-orientiert-richard-diebel-nils-berns-andreas-christ-wolfgang-meier-magdalena-turska-lars-windauer}}
\sphinxAtStartPar
Im Rahmen eines Workshops zur neuen Version des TEI Publishers konnten wir das Tool direkt erproben. Der TEI Publisher ist eine auf eXist\sphinxhyphen{}db basierende Open\sphinxhyphen{}Source\sphinxhyphen{}Plattform, die es ermöglicht, TEI\sphinxhyphen{}XML\sphinxhyphen{}Dokumente ohne tiefgreifende Webentwicklungskenntnisse in ansprechende digitale Editionen zu überführen. Die neue Version setzt auf einen modularen Ansatz mit konfigurierbaren Profilen, die sowohl die Wartung vereinfachen als auch die Einstiegshürde für eigene Projekte deutlich senken. In der hands\sphinxhyphen{}on Session wurde schnell deutlich, wie niedrigschwellig sich mittlerweile visuell ansprechende Editionsanwendungen erstellen lassen.


\section{Doctoral Consortium: Möglichkeiten und Potenziale von inter\sphinxhyphen{} und transmedialen Editionen (Dennis Friedl)}
\label{\detokenize{tagung/digitale-editionen:doctoral-consortium-moglichkeiten-und-potenziale-von-inter-und-transmedialen-editionen-dennis-friedl}}
\sphinxAtStartPar
Am Beispiel der \sphinxstyleemphasis{Erich Wolfgang Korngold}\sphinxhyphen{}Werkausgabe diskutierte der Vortrag das Konzept der hybriden Edition, bei der digitale und gedruckte Varianten nicht bloß parallel existieren, sondern gezielt zusammenspielen – etwa über QR\sphinxhyphen{}Codes, die beide Medien miteinander verknüpfen. Interessant ist hier auch die Überlegung, die spezifischen Eigenarten einzelner Ausgabemedien produktiv zu nutzen: Die Mobilität und Touch\sphinxhyphen{}Interaktion des Smartphones, die Vertikalität des Displays oder die Immersion von VR/AR eröffnen jeweils eigene Formen des Editionskonsums.

\sphinxstepscope


\chapter{Andere Highlights}
\label{\detokenize{tagung/anderes:andere-highlights}}\label{\detokenize{tagung/anderes::doc}}
\sphinxAtStartPar
Neben den thematischen Schwerpunkten, die ich mir gesetzt habe, gab es auch viele weitere spannende Beiträge, die ich nicht unerwähnt lassen möchte. Gerade bei Vorträgen, bei denen ich eher zufällig war (beispielsweise, weil andere Sitzungen schon komplett “ausgebucht” waren), habe ich viele Impulse für mich mitgenommen – ganz im Sinne der Suchstrategie des Browsings, die zu Serependität führt (wie das Anna Schlander in ihrem Vortrag wunderbar ausgeführt hat).


\section{Keynote: Data As \_\_\_\_\_\_\_: Exploring the Plurality of Data in Visualization (Miriah Meyer)}
\label{\detokenize{tagung/anderes:keynote-data-as-exploring-the-plurality-of-data-in-visualization-miriah-meyer}}
\sphinxAtStartPar
In ihrer Keynote plädierte Miriah Meyer dafür, Daten nicht (nur) als neutrale, objektive Grundlage zu verstehen, sondern als vielschichtiges, mit Geschichte, Technologie und Menschen verwobenes Material, das Macht ausübt und immer eine Perspektive trägt. Vorgestellt wurde dabei insbesondere das didaktisch interessante Konzept des Data Crafting – also das bewusste, handwerkliche Gestalten und Visualisieren von Daten als Reflexionsprozess, der Datenkompetenz fördern kann.


\section{Doctoral Consortium: Die Vergangenheit umgibt uns – Digitalisiertes Kulturerbe durch situierte Visualisierung vor Ort erfahrbar machen (Markus Passecker)}
\label{\detokenize{tagung/anderes:doctoral-consortium-die-vergangenheit-umgibt-uns-digitalisiertes-kulturerbe-durch-situierte-visualisierung-vor-ort-erfahrbar-machen-markus-passecker}}
\sphinxAtStartPar
Markus Passecker beschäftigt sich mit der Frage, wie die großen Mengen digitalisierten Kulturerbes, die heute in Online\sphinxhyphen{}Sammlungen verfügbar sind, wieder in ihren ursprünglichen räumlichen Kontext zurückgebracht werden können. Sein Ansatz setzt auf situierte Visualisierung und Mixed Reality: Digitale Objekte werden nicht im Museum oder am Bildschirm betrachtet, sondern an relevantem Ort erfahrbar gemacht – von der gezielten Erkundung bis zur zufälligen Entdeckung. Am Beispiel des Projekts \sphinxstyleemphasis{Mapping Austrofascism} wurden dabei praktische Herausforderungen greifbar, etwa die Reduktion von Visual Overload durch Filtering und räumliche Verschiebungen oder die Erkennung sicherer Zonen für die Overlay\sphinxhyphen{}Platzierung. Er stellt die Frage nach einer angemessenen Evaluierungsstrategie – ob also Nutzerengagement, emotionale Verbundenheit oder Erinnerungsleistung der geeignetere Maßstab für solche immersiven Kulturerbe\sphinxhyphen{}Erfahrungen ist.


\section{Zwischen Struktur und Überraschung – Gestaltung explorativer Suchfunktionen für digitale Korpora (Anna Schlander)}
\label{\detokenize{tagung/anderes:zwischen-struktur-und-uberraschung-gestaltung-explorativer-suchfunktionen-fur-digitale-korpora-anna-schlander}}
\sphinxAtStartPar
Anna Schlander näherte sich der Frage, wie explorative Suche in digitalen Korpora gestaltet werden kann, vom Konzept des Browsings aus – verstanden nicht als zielloses Stöbern, sondern als eine Form der Recherche, bei der Suchkriterien nur teilweise bekannt sind und Neugier und Serendipität bewusst als Motivation einbezogen werden. Am Beispiel des PubMed Cliques Explorers zeigte sie, wie sich solche Prinzipien technisch umsetzen lassen: Ein graphbasierter Workflow – von CSV über RDF bis zur algorithmischen Detektion dichter Subgraphen – macht es möglich, thematische Cluster zu erkunden.


\section{Warum wir in den Digital Humanities messen (sollten) (Evelyn Gius)}
\label{\detokenize{tagung/anderes:warum-wir-in-den-digital-humanities-messen-sollten-evelyn-gius}}
\sphinxAtStartPar
Evelyn Gius argumentierte, dass Messen in den Digital Humanities – verstanden als die Transformation empirischer Eigenschaften in Information – eine sinnvollere Rahmung bietet als der klassische Begriff der Operationalisierung, da er Unsicherheiten explizit mitdenkt. Besonders interessant für mich war hier das Beispiel von Sentiment\sphinxhyphen{}Analyse und LLMs, wo vielfältige Messunsicherheiten (bezüglich Instrument, Kalibrierung, Interaktion und Skalierung) lauern. Deren bewusste Reflexion sollte jedoch nicht als Problem, sondern als strukturierendes Prinzip für den geisteswissenschaftlichen Analyseworkflow verstanden werden.


\section{The Dialogue between Reason and Intuition in Contemporary Philosophy (Maximilian Noichl, Simon DeDeo)}
\label{\detokenize{tagung/anderes:the-dialogue-between-reason-and-intuition-in-contemporary-philosophy-maximilian-noichl-simon-dedeo}}
\sphinxAtStartPar
Maximilian Noichl und Simon DeDeo untersuchten, wie zwei grundlegende Modi philosophischen Arbeitens – das Aufstellen von Argumenten und Positionen einerseits (\sphinxstyleemphasis{reason}), das Entwickeln von Intuitionen durch Metaphern und Gedankenexperimente andererseits (\sphinxstyleemphasis{intuition}) – in philosophischen Texten zusammenspielen und welcher der beiden stabiler ist. Methodisch interessant war die Kombination aus GPT\sphinxhyphen{}4o\sphinxhyphen{}basierter Annotation, einem eigens trainierten ModernBERT\sphinxhyphen{}Modell und KNN\sphinxhyphen{}Graphen zur Zusammenführung referenzierter Positionen und Beispiele über einen Korpus von JSTOR\sphinxhyphen{}Philosophietexten ab 1950 – ein Workflow, der zeigt, wie leistungsfähig hybride Ansätze aus klassischem NLP und generativen Modellen für geisteswissenschaftliche Fragestellungen sein können.


\section{Projektmanagement als Methode (Julia Alili)}
\label{\detokenize{tagung/anderes:projektmanagement-als-methode-julia-alili}}
\sphinxAtStartPar
Julia Alili argumentierte, dass Projektmanagement in den Digital Humanities nicht bloß organisatorisches Beiwerk ist, sondern als genuine wissenschaftliche Methode verstanden werden sollte. Ausgehend von typischen Herausforderungen in DH\sphinxhyphen{}Projekten – Interdisziplinarität, uneinheitliche Begrifflichkeiten, lückenhafte Dokumentation – plädierte sie für eine Adaption agiler Ansätze, bei der Fortschritt nicht in funktionierender Software gemessen wird, sondern in neuen Erkenntnissen, Standards und Dokumenten. Wichtig ist hier die Idee eines kontinuierlichen Wissensspeichers durch Requirements Engineering: Glossare, Entscheidungsbegründungen und Versionierungen schaffen Transparenz und machen Projektentscheidungen nachvollziehbar. Diese agile Form von DH\sphinxhyphen{}Projektmanagement möchte sie konsequenterweise auch in der DH\sphinxhyphen{}Ausbildung verankert sehen.

\sphinxstepscope


\chapter{Fazit zur Tagung}
\label{\detokenize{tagung/fazit:fazit-zur-tagung}}\label{\detokenize{tagung/fazit::doc}}
\sphinxAtStartPar
Die DHd 2026 hat einmal mehr gezeigt, wie lebendig und vielseitig das Feld der Digital Humanities ist. Die thematische Bandbreite der Beiträge war beeindruckend, und ich nehme viele Impulse mit – für Forschung ebenso wie für die Lehre.

\sphinxAtStartPar
Nicht weniger wertvoll war aber das Drumherum: Die Organisation und das Ambiente haben eine tolle Atmosphäre geschaffen. Besonders schön war die Möglichkeit zur Vernetzung innerhalb der Community, insbesondere auch im Rahmen des Stipendienprogramms.

\sphinxAtStartPar
Nicht zuletzt empfand ich es als extrem bereichernd, eigene Beiträge vorzustellen – sowohl innerhalb der Postersession (\sphinxhref{https://doi.org/10.5281/zenodo.18999645}{DH on the Edge?}) als auch im Rahmen der Session zu Historical Perspectives (\sphinxhref{https://doi.org/10.5281/zenodo.18696271}{“I was JUST Cleopatra man” – ein quantitativer Zugang zu KI\sphinxhyphen{}generierten Geschichtsnarrativen und ihrer Rezeption auf TikTok}). Die Fragen und Rückmeldungen aus der Community haben mir neue Perspektiven auf meine Arbeit eröffnet, und ich habe die DH\sphinxhyphen{}Community einmal mehr als ausgesprochen offen, neugierig und diskussionsfreudig erlebt.

\begin{sphinxuseclass}{sd-tab-set}
\begin{sphinxuseclass}{sd-tab-item}\subsubsection*{Universität Wien}

\begin{sphinxuseclass}{sd-tab-content}
\sphinxAtStartPar
\sphinxincludegraphics{{1}.jpg}

\end{sphinxuseclass}
\end{sphinxuseclass}
\begin{sphinxuseclass}{sd-tab-item}\subsubsection*{Festsaal}

\begin{sphinxuseclass}{sd-tab-content}
\sphinxAtStartPar
\sphinxincludegraphics{{2}.jpg}

\end{sphinxuseclass}
\end{sphinxuseclass}
\begin{sphinxuseclass}{sd-tab-item}\subsubsection*{Empfang im Rathaus}

\begin{sphinxuseclass}{sd-tab-content}
\sphinxAtStartPar
\sphinxincludegraphics{{3}.jpg}

\end{sphinxuseclass}
\end{sphinxuseclass}
\end{sphinxuseclass}
\sphinxstepscope


\part{Das Mentoring\sphinxhyphen{}Programm}

\sphinxstepscope


\chapter{Das DHd\sphinxhyphen{}Mentoringprogramm}
\label{\detokenize{mentoring/programm:das-dhd-mentoringprogramm}}\label{\detokenize{mentoring/programm::doc}}
\sphinxAtStartPar
Der DHd\sphinxhyphen{}Verband schreibt jährlich ein \sphinxhref{https://digitalhumanities.de/mentoring/}{Mentoringprogramm} für Early Career Wissenschaftler:innen in den Digital Humanities aus. Ziel des Programms ist es, Nachwuchswissenschaftler:innen bei der Orientierung und Positionierung in einem interdisziplinären Forschungsfeld zu unterstützen, in dem sich fachdisziplinäre und digitale Perspektiven überschneiden.

\sphinxAtStartPar
Ich habe am Mentoringprogramm im Jahrgang 2025/2026 teilgenommen.


\section{An wen richtet sich das Programm?}
\label{\detokenize{mentoring/programm:an-wen-richtet-sich-das-programm}}
\sphinxAtStartPar
Das Mentoringprogramm richtet sich insbesondere an:
\begin{itemize}
\item {} 
\sphinxAtStartPar
Promovierende in den Digital Humanities

\item {} 
\sphinxAtStartPar
fortgeschrittene Masterstudierende

\item {} 
\sphinxAtStartPar
PostDocs

\item {} 
\sphinxAtStartPar
Wissenschaftler:innen, die digitale Methoden einsetzen und bislang wenig Anbindung an die DH\sphinxhyphen{}Community haben

\end{itemize}


\section{Wobei unterstützt das Mentoring?}
\label{\detokenize{mentoring/programm:wobei-unterstutzt-das-mentoring}}
\sphinxAtStartPar
Das Programm bietet Unterstützung bei:
\begin{itemize}
\item {} 
\sphinxAtStartPar
der fachlichen Profilbildung an der Schnittstelle von Disziplin und Digital Humanities

\item {} 
\sphinxAtStartPar
der Karriereplanung innerhalb und außerhalb der Wissenschaft

\item {} 
\sphinxAtStartPar
Fragen zu Methoden, Tools und Forschungsinfrastrukturen

\item {} 
\sphinxAtStartPar
der Orientierung in der DH\sphinxhyphen{}Community und dem Aufbau von Netzwerken

\end{itemize}


\section{Formate und Veranstaltungen}
\label{\detokenize{mentoring/programm:formate-und-veranstaltungen}}
\sphinxAtStartPar
Das Mentoringprogramm kombiniert mehrere Formate:


\subsection{Individuelles Mentoring}
\label{\detokenize{mentoring/programm:individuelles-mentoring}}\begin{itemize}
\item {} 
\sphinxAtStartPar
Begleitung durch eine:n Mentor:in über 12 Monate

\item {} 
\sphinxAtStartPar
gemeinsame Festlegung von Zielen zu Beginn des Mentorings

\item {} 
\sphinxAtStartPar
mindestens vier strukturierte (Online\sphinxhyphen{})Treffen

\item {} 
\sphinxAtStartPar
thematische Ausrichtung nach den individuellen Bedürfnissen der Mentees

\end{itemize}


\subsection{Peer\sphinxhyphen{}to\sphinxhyphen{}Peer\sphinxhyphen{}Mentoring}
\label{\detokenize{mentoring/programm:peer-to-peer-mentoring}}\begin{itemize}
\item {} 
\sphinxAtStartPar
Vernetzung der Mentees untereinander

\item {} 
\sphinxAtStartPar
Austausch von Erfahrungen und Ressourcen

\item {} 
\sphinxAtStartPar
informelle Zusammenarbeit zu Tools, Methoden oder Fragestellungen

\end{itemize}


\subsection{Begleitende Veranstaltungen}
\label{\detokenize{mentoring/programm:begleitende-veranstaltungen}}\begin{itemize}
\item {} 
\sphinxAtStartPar
Kick\sphinxhyphen{}off\sphinxhyphen{}Veranstaltung mit einer Mentoring\sphinxhyphen{}Expertin zu Beginn des Mentoringjahres

\item {} 
\sphinxAtStartPar
Meet and Greet im Rahmen der DHd\sphinxhyphen{}Jahrestagung

\end{itemize}


\section{Was ist Mentoring?}
\label{\detokenize{mentoring/programm:was-ist-mentoring}}
\sphinxAtStartPar
Mentoring ist eine unterstützende Beziehung zwischen einer erfahrenen Person (Mentor:in) und einer weniger erfahrenen Person (Mentee), die auf freiwilligem Austausch, Vertrauen und Reflexion basiert. Mentor:innen fungieren dabei als Vorbilder, geben Feedback und teilen ihre eigenen Erfahrungen im Feld der Digital Humanities.

\sphinxAtStartPar
Ziel des Mentorings ist es, Mentees dabei zu unterstützen,
\begin{itemize}
\item {} 
\sphinxAtStartPar
ihre Karriereziele und \sphinxhyphen{}optionen zu reflektieren,

\item {} 
\sphinxAtStartPar
Prioritäten zu setzen und persönliche Entwicklungsfelder zu identifizieren,

\item {} 
\sphinxAtStartPar
eine Richtung für die weitere Karriereentwicklung zu entwickeln,

\item {} 
\sphinxAtStartPar
Motivation, Orientierung und konstruktives Feedback zu erhalten.

\end{itemize}

\sphinxAtStartPar
Ein zentraler Bestandteil des Mentorings ist die Weitergabe von informellem Wissen (\sphinxstyleemphasis{tacit knowledge})
über Strukturen, Erwartungen und Abläufe im wissenschaftlichen und DH\sphinxhyphen{}spezifischen Umfeld, das in formalen Kontexten oft nicht explizit vermittelt wird. Mentoring versteht sich dabei als begleitender Coachingprozess auf Augenhöhe – nicht als Bewertung oder Kontrolle.


\section{Was ist Mentoring nicht?}
\label{\detokenize{mentoring/programm:was-ist-mentoring-nicht}}
\sphinxAtStartPar
Mentoring ist keine Betreuung im formalen Sinne. Mentor:innen übernehmen keine fachliche oder institutionelle Verantwortung für die Arbeit der Mentees, etwa in Form einer Co\sphinxhyphen{}Betreuung von Abschlussarbeiten oder Dissertationen. Eine solche Konstellation würde einen Rollenkonflikt darstellen und ist ausdrücklich nicht Ziel des Mentorings.

\sphinxAtStartPar
Ebenso ist Mentoring nicht gleichzusetzen mit Sponsorship. Während Sponsor:innen aktiv Karrierechancen eröffnen – etwa durch Empfehlungsschreiben, gezielte Förderungen oder den Aufbau konkreter Kooperationen –, bleiben Mentor:innen primär in einer beratenden und reflektierenden Rolle. Eine Ausweitung in Richtung Sponsorship kann sich im Einzelfall ergeben, ist jedoch weder vorgesehen noch verpflichtend.

\sphinxAtStartPar
Darüber hinaus sind Mentor:innen
\begin{itemize}
\item {} 
\sphinxAtStartPar
keine Mediator:innen in Konfliktsituationen,

\item {} 
\sphinxAtStartPar
keine psychologischen Berater:innen,

\item {} 
\sphinxAtStartPar
und nicht zuständig für die Lösung persönlicher oder institutioneller Konflikte.

\end{itemize}

\sphinxAtStartPar
Sie können jedoch dabei helfen, Situationen einzuordnen, Bewältigungsstrategien zu reflektieren oder auf geeignete Anlaufstellen hinzuweisen. Für sehr allgemeine oder nicht mentorspezifische Fragen bietet sich zudem der Austausch innerhalb der Peer\sphinxhyphen{}Group an.

\sphinxAtStartPar
Ein grundlegendes Prinzip des Mentorings ist die Vertraulichkeit: Inhalte aus Mentoringgesprächen werden nicht
nach außen getragen und bleiben innerhalb der Mentoringbeziehung.


\section{Weiterführende Links}
\label{\detokenize{mentoring/programm:weiterfuhrende-links}}\begin{itemize}
\item {} 
\sphinxAtStartPar
\sphinxurl{https://digitalhumanities.de/mentoring/}

\item {} 
\sphinxAtStartPar
Ausschreibung 2025/2026: \sphinxurl{https://digitalhumanities.de/2025/04/04/ausschreibung-zum-dhd-early-career-mentoring-programm-2025-26/}

\end{itemize}

\sphinxstepscope


\chapter{Meine Erfahrungen}
\label{\detokenize{mentoring/meine-erfahrungen:meine-erfahrungen}}\label{\detokenize{mentoring/meine-erfahrungen::doc}}

\section{Kick\sphinxhyphen{}Off}
\label{\detokenize{mentoring/meine-erfahrungen:kick-off}}
\sphinxAtStartPar
Im September startete das Programm mit einem Coaching von Dr. Claudia Eilles\sphinxhyphen{}Matthiessen. Dadurch wurden wir von Anfang an dafür sensibiliert, was ein erfolgreiches Mentoring ausmacht, und erhielten zahlreiche praktische Tipps und Tools. Gleichzeitig bot der Kick\sphinxhyphen{}off\sphinxhyphen{}Termin eine gute Gelegenheit, die anderen Mentor*innen und Mentees kennenzulernen.


\section{Meetings}
\label{\detokenize{mentoring/meine-erfahrungen:meetings}}
\sphinxAtStartPar
Mit \sphinxhref{https://cceh.uni-koeln.de/personen/elisa-cugliana/}{Junior\sphinxhyphen{}Professorin Dr. Elisa Cugliana} wurde eine Mentorin für mich gefunden, die fachlich ideal passt – ihre Schwerpunkte liegen auf der Linguistik und Philologie germanischer Sprachen und der digitalen Editionswissenschaft.
Gemeinsam haben Elisa und ich festgelegt, dass mich das Mentoring primär bei Themenfindung für und Start in die Promotionsphase begleiten soll. Dazu sprechen wir circa alle 6 Wochen.
Ich empfinde diese Gespräche, nicht zuletzt wegen Elisas offener und unterstützender Art, als sehr bereichernd und ungezwungen. Es war extrem hilfreich, eine externe Sicht und Beurteilung der Dinge, die trotzdem aus der Disziplin kommt, zu erfahren.
Teil des Programms ist es auch, diese Gespräche (zumindest grob) zu dokumentieren. Auch das empfinde ich als sehr wertvoll – nicht nur zum Strukturieren und Festhalten von Gedanken, sondern auch zur Vergegenwärtigung von Fortschritt.


\section{Fazit}
\label{\detokenize{mentoring/meine-erfahrungen:fazit}}
\sphinxAtStartPar
Das Mentoring\sphinxhyphen{}Programm war für mich insgesamt eine sehr bereichernde Erfahrung. Elisa hat mich mit großem Engagement begleitet und wertvolle Impulse gegeben. Stellenweise hätte ich mir noch mehr Raum für den Peer\sphinxhyphen{}to\sphinxhyphen{}Peer\sphinxhyphen{}Austausch unter den Menteees gewünscht, um voneinander zu lernen und Erfahrungen zu teilen – insgesamt bin ich aber sehr froh und dankbar, Teil eines so tollen Programms zu sein.

\sphinxstepscope


\part{Arbeitsgruppen des DHd\sphinxhyphen{}Verbandes}

\sphinxstepscope


\chapter{Übersicht zu den DHd\sphinxhyphen{}Arbeitsgruppen}
\label{\detokenize{arbeitsgruppen/uebersicht:ubersicht-zu-den-dhd-arbeitsgruppen}}\label{\detokenize{arbeitsgruppen/uebersicht::doc}}

\begin{savenotes}\sphinxattablestart
\sphinxthistablewithglobalstyle
\centering
\begin{tabulary}{\linewidth}[t]{TTT}
\sphinxtoprule
\sphinxstyletheadfamily 
\sphinxAtStartPar
Name
&\sphinxstyletheadfamily 
\sphinxAtStartPar
Ziele
&\sphinxstyletheadfamily 
\sphinxAtStartPar
Link
\\
\sphinxmidrule
\sphinxtableatstartofbodyhook
\sphinxAtStartPar
Angewandte Generative KI in den Digitalen Geisteswissenschaften (AGKI\sphinxhyphen{}DH)
&
\sphinxAtStartPar
\sphinxhyphen{} Identifikation, Bewertung und Kategorisierung des Anwendungsspektrums von generativer KI in den DH (Analyse, Modellierung, Datengenerierung etc.)\sphinxhyphen{} Erarbeitung von Best Practices im Umgang mit generativer KI\sphinxhyphen{} Entwicklung von Evaluationskriterien und \sphinxhyphen{}szenarien für unterschiedliche Aufgaben\sphinxhyphen{} Diskussion und Austausch von Erfahrungen, Experimenten und Fehlschlägen\sphinxhyphen{} Wissenstransfer: Training und Lehre für die Community\sphinxhyphen{} Abwägung von Chancen und Risiken\sphinxhyphen{} Klärung und Erklärung von Begriffen für die DH Community
&
\sphinxAtStartPar
\sphinxurl{https://agki-dh.github.io/}
\\
\sphinxhline
\sphinxAtStartPar
Digital Humanities Theorie
&
\sphinxAtStartPar
\sphinxhyphen{} Schaffung eines strukturierten Kommunikationsrahmens für den interdisziplinären Austausch über Theoriebildung in den DH\sphinxhyphen{} Bestandsaufnahme, Bewertung und Sichtbarmachung vielfältiger theoretischer Zugänge in den DH\sphinxhyphen{} Wissenschaftstheoretische und \sphinxhyphen{}soziologische Reflexion des Verhältnisses der DH zu Nachbardisziplinen (Informatik, Medienwissenschaften u. a.)\sphinxhyphen{} Auseinandersetzung mit kollektiven und partizipativen Formen der Theoriearbeit\sphinxhyphen{} Dissemination der Ergebnisse in verschiedenen Formaten (Paper, Blogbeiträge) über virtuelle und persönliche Kollaborationsformate
&
\sphinxAtStartPar
\sphinxurl{https://dhtheorien.hypotheses.org/}
\\
\sphinxhline
\sphinxAtStartPar
Digitales Museum
&
\sphinxAtStartPar
\sphinxhyphen{} Auseinandersetzung mit den unterschiedlichen Rollen von Museen in den DH\sphinxhyphen{} Herausarbeitung der Bedeutung und Potenziale von Museen als DH\sphinxhyphen{}Akteure\sphinxhyphen{} Diskussion der Transformation zur Institution „Digitales Museum”\sphinxhyphen{} Funktion als Vermittler und Beratungsgremium innerhalb der DH\sphinxhyphen{}Community
&
\sphinxAtStartPar
\sphinxurl{https://dig-hum.de/ag-digitales-museum}
\\
\sphinxhline
\sphinxAtStartPar
Digitales Publizieren
&
\sphinxAtStartPar
\sphinxhyphen{} Sichtbarmachung des Mehrwerts digitaler Publikationen\sphinxhyphen{} Erarbeitung von Best Practices (FAIR\sphinxhyphen{}Prinzipien, Langzeitarchivierung) für digitale Publikationen\sphinxhyphen{} Reflexion des Wandels wissenschaftlicher Autorschaft durch kollaborative Textproduktion, Annotation und automatisierte Verfahren\sphinxhyphen{} Förderung von Qualität, Nachnutzbarkeit, Reputation und Langzeitverfügbarkeit digitaler Publikationen\sphinxhyphen{} Auseinandersetzung mit ethischen und ökonomischen Fragen, insbesondere Open Access\sphinxhyphen{} Entwicklung von Handlungsempfehlungen für die Community sowie Empfehlungen an Förderer und Institutionen
&
\sphinxAtStartPar
\sphinxurl{https://dig-hum.de/ag-digitales-publizieren}
\\
\sphinxhline
\sphinxAtStartPar
Digitale Wissenschaftskommunikation und Public Humanities (AG\sphinxhyphen{}WissKomm)
&
\sphinxAtStartPar
\sphinxhyphen{} Bündelung von Aktivitäten der deutschsprachigen DH\sphinxhyphen{}Community im Bereich Wissenschaftskommunikation und Public Humanities\sphinxhyphen{} Auswertung und Reflexion von Wissenschaftskommunikation und Public Humanities innerhalb der DH\sphinxhyphen{} Stärkung der öffentlichen Wahrnehmung und Relevanz der Digital Humanities\sphinxhyphen{} Entwicklung von Best Practices für den Einstieg in die Wissenschaftskommunikation\sphinxhyphen{} Reflexion des Mehrwerts von Wissenschaftskommunikation für Forschende und Projekte\sphinxhyphen{} Einbeziehung der Forschungs\sphinxhyphen{} und allgemeinen Öffentlichkeit in DH\sphinxhyphen{}Projekte (Citizen Science / Public Digital Humanities)
&
\sphinxAtStartPar
\sphinxurl{https://publicdh.hypotheses.org/}
\\
\sphinxhline
\sphinxAtStartPar
Digitale 3D\sphinxhyphen{}Rekonstruktion
&
\sphinxAtStartPar
\sphinxhyphen{} Vernetzung der Akteure im deutschsprachigen Raum im Bereich digitale 3D\sphinxhyphen{}Rekonstruktion von Kulturerbe\sphinxhyphen{} Begriffsklärung und Entwicklung einer gemeinsamen Arbeitsmethodologie\sphinxhyphen{} Erarbeitung von Dokumentationsstandards und Lösungen zur Langzeitarchivierung von Rekonstruktionsprojekten und 3D\sphinxhyphen{}Datensätzen\sphinxhyphen{} Förderung gemeinsamer Publikationen und Auftritte
&
\sphinxAtStartPar

\\
\sphinxhline
\sphinxAtStartPar
Empowerment
&
\sphinxAtStartPar
\sphinxhyphen{} Schaffung eines Gesprächsraums für die Auseinandersetzung mit Machtstrukturen, prekären Arbeitsverhältnissen und diskriminatorischen Praktiken in der deutschsprachigen DH\sphinxhyphen{}Community\sphinxhyphen{} Reflexion von Geschlechtergerechtigkeit, kolonialem Erbe und weiteren wissenschaftspolitischen Fragen im DH\sphinxhyphen{}Kontext\sphinxhyphen{} Anschluss an internationale Diskurse zu Critical Digital Humanities\sphinxhyphen{} Förderung von Empowerment innerhalb der Community – gemeinsam, konstruktiv und pragmatisch
&
\sphinxAtStartPar
\sphinxurl{https://empowerdh.github.io/}
\\
\sphinxhline
\sphinxAtStartPar
Film und Video
&
\sphinxAtStartPar
\sphinxhyphen{} Vernetzung von Projekten und Akteuren, die computergestützte Verfahren in der Filmwissenschaft und der Arbeit mit Video\sphinxhyphen{}Ressourcen einsetzen\sphinxhyphen{} Zusammentragen und Systematisierung bestehender Ansätze computergestützter Filmanalyse im deutschsprachigen Raum\sphinxhyphen{} Identifikation zentraler Problembereiche und Herausforderungen (Ressourcenbedarf, mediale Komplexität, Toolentwicklung)\sphinxhyphen{} Öffentliche Zugänglichmachung und Verbreitung von Ergebnissen\sphinxhyphen{} Anstoßen und Unterstützen neuer Forschungsaktivitäten im Bereich Digital Film Studies
&
\sphinxAtStartPar
\sphinxurl{https://dhdagfilm.hypotheses.org/}
\\
\sphinxhline
\sphinxAtStartPar
Geisteswissenschaftliches Forschungsdatenmanagement (gwFDM)
&
\sphinxAtStartPar
\sphinxhyphen{} Vernetzung von Personen, Einrichtungen und Infrastrukturen im Bereich geisteswissenschaftliches Forschungsdatenmanagement (FDM)\sphinxhyphen{} Aufbereitung und Bereitstellung von Forschungsdaten gemäß FAIR\sphinxhyphen{}Prinzipien\sphinxhyphen{} Förderung des Austauschs über Begriffe, Selbstverständnisse und Problemlagen im FDM\sphinxhyphen{} Identifikation von Lücken in der FDM\sphinxhyphen{}Versorgungslandschaft\sphinxhyphen{} Organisation der Tagungsreihe FORGE – Forschungsdaten in den Geisteswissenschaften
&
\sphinxAtStartPar
\sphinxurl{https://dhd-ag-datenzentren.github.io/}
\\
\sphinxhline
\sphinxAtStartPar
Graphen \& Netzwerke
&
\sphinxAtStartPar
\sphinxhyphen{} Förderung des Austauschs zu Entwicklung und Anwendung von Graphen und Netzwerken in den DH\sphinxhyphen{} Erarbeitung von Modellierungsbeispielen und Standards für graphbasierte Wissensmodellierung\sphinxhyphen{} Reflexion des Mehrwerts von Graphen für die Textmodellierung (mehrdimensionale Hierarchien, semantische Erschließung)\sphinxhyphen{} Auseinandersetzung mit Versionierung, Zitierfähigkeit und Langzeitarchivierung graphbasierter Daten\sphinxhyphen{} Diskussion von Analysepotentialen durch Graphen und Netzwerkmetriken
&
\sphinxAtStartPar
\sphinxurl{https://graphentechnologien.hypotheses.org/}
\\
\sphinxhline
\sphinxAtStartPar
Greening DH
&
\sphinxAtStartPar
\sphinxhyphen{} Evaluation bestehender Rechenmethoden, Tools und Maßnahmen im Hinblick auf ökologische Nachhaltigkeit\sphinxhyphen{} Erarbeitung fachspezifischer Best\sphinxhyphen{}Practice\sphinxhyphen{}Richtlinien für unterschiedliche Tätigkeitsfelder (Reisen, Konferenzen, Hosting, Materialbeschaffung)\sphinxhyphen{} Förderung des Austauschs mit Wissenschaftler*innen anderer Fächer sowie internationaler Interessengruppen\sphinxhyphen{} Positionierung der DH als Vorreiterin für nachhaltigere Forschungspraktiken
&
\sphinxAtStartPar
\sphinxurl{https://dhd-greening.github.io/}
\\
\sphinxhline
\sphinxAtStartPar
Multilingual DH
&
\sphinxAtStartPar
\sphinxhyphen{} Kritische Auseinandersetzung mit der Zentrierung der DH auf lateinbasierte Schriftsprachen und englischsprachige Infrastrukturen\sphinxhyphen{} Thematisierung kolonialer Strukturen und der Marginalisierung des globalen Südens in der DH\sphinxhyphen{}Wissenschaftspraxis\sphinxhyphen{} Förderung der Vernetzung multilingualer und nicht\sphinxhyphen{}lateinschriftlicher DH\sphinxhyphen{}Projekte zur Vermeidung von Parallelentwicklungen\sphinxhyphen{} Anschluss an internationale Diskurse zu “Decolonizing DH” und “Disrupting Digital Monolingualism”\sphinxhyphen{} Entwicklung von Ideen und Maßnahmen zur Verbesserung der Infrastruktur und Repräsentation nicht\sphinxhyphen{}lateinschriftlicher Sprachen
&
\sphinxAtStartPar
\sphinxurl{http://multilingualdh.de}
\\
\sphinxhline
\sphinxAtStartPar
OCR
&
\sphinxAtStartPar
\sphinxhyphen{} Vernetzung von Nutzer*innen, Vermittler*innen, Entwickler*innen und Forscher*innen im Bereich OCR\sphinxhyphen{} Erarbeitung und Verbreitung von Best Practices zu OCR\sphinxhyphen{}Workflows, Formaten und Parametern\sphinxhyphen{} Technologietransfer aus der Entwicklung in die Praxis\sphinxhyphen{} Vermittlung zwischen geisteswissenschaftlichen Anforderungen und technischen Möglichkeiten der OCR\sphinxhyphen{}Softwareentwicklung\sphinxhyphen{} Kollaborative Bearbeitung offener Forschungsfragen in Zusammenarbeit mit verwandten AGs (z. B. AG Zeitungen \& Zeitschriften)\sphinxhyphen{} Fokus auf freier und quelloffener Software und Publikationen
&
\sphinxAtStartPar
\sphinxurl{https://dhd-ag-ocr.github.io}
\\
\sphinxhline
\sphinxAtStartPar
Referenzcurriculum Digital Humanities
&
\sphinxAtStartPar
\sphinxhyphen{} Entwicklung eines Referenzcurriculums als gemeinsames Orientierungsmodell für DH\sphinxhyphen{}Studiengänge im deutschsprachigen Raum\sphinxhyphen{} Erleichterung der wechselseitigen Anerkennung von Studienleistungen durch studiengangsunabhängige Definitionen\sphinxhyphen{} Unterstützung der Akkreditierung von Studiengängen durch institutionsübergreifende Referenzwerte\sphinxhyphen{} Diskussion aktueller Fragen rund um Studium und Lehre in den DH\sphinxhyphen{} Proaktive Selbstdefinition der DH\sphinxhyphen{}Ausbildung, um externe Fremdbestimmung zu vermeiden
&
\sphinxAtStartPar
\sphinxurl{http://dig-hum.de/ag-referenzcurriculum-digital-humanities}
\\
\sphinxhline
\sphinxAtStartPar
Research Software Engineering in den Digital Humanities (DH\sphinxhyphen{}RSE)
&
\sphinxAtStartPar
\sphinxhyphen{} Förderung qualitativ hochwertiger und nachhaltiger Softwareentwicklung in den DH\sphinxhyphen{} Stärkung der akademischen Position von Softwareentwickler*innen in den Geistes\sphinxhyphen{} und Kulturwissenschaften\sphinxhyphen{} Vernetzung mit der nationalen und internationalen Research Software Engineers (RSE)\sphinxhyphen{}Community
&
\sphinxAtStartPar
\sphinxurl{https://dh-rse.github.io/}
\\
\sphinxhline
\sphinxAtStartPar
Spiele
&
\sphinxAtStartPar
\sphinxhyphen{} Erfassung von Computer\sphinxhyphen{}, Video\sphinxhyphen{} und elektronischen Spielen als Gegenstand der Digital Humanities\sphinxhyphen{} Entwicklung medienadäquater Methoden und Werkzeuge unter Berücksichtigung von Interaktivität, Operativität und technischer Unterfläche\sphinxhyphen{} Förderung einer umfassenden Cultural Heritage Preservation von Spielen inkl. Spielkulturen und paratextuellen Prozessen\sphinxhyphen{} Interdisziplinäre Kollaboration zwischen technischen, kultur\sphinxhyphen{} und geisteswissenschaftlichen Disziplinen sowie Museen und Archiven\sphinxhyphen{} Publikation von Ergebnissen in Büchern, Zeitschriften und einem AG\sphinxhyphen{}Blog sowie Etablierung einer regelmäßigen Workshopreihe
&
\sphinxAtStartPar
\sphinxurl{http://dhspiele.de}
\\
\sphinxhline
\sphinxAtStartPar
Zeitungen \& Zeitschriften
&
\sphinxAtStartPar
\sphinxhyphen{} Vernetzung von Institutionen (Bibliotheken), Projekten und Einzelforschenden im Bereich Zeitungs\sphinxhyphen{} und Zeitschriftenforschung\sphinxhyphen{} Vergleich und Vermittlung von Best\sphinxhyphen{}Practice\sphinxhyphen{}Lösungen aus bestehenden Projekten\sphinxhyphen{} Erarbeitung von Standards und Entwicklungsempfehlungen für die digitale Zeitungs\sphinxhyphen{} und Zeitschriftenforschung\sphinxhyphen{} Förderung nachhaltiger Speicherung und Zugänglichkeit von Forschungsdaten auch über die DH\sphinxhyphen{}Community hinaus
&
\sphinxAtStartPar
\sphinxurl{https://dhd-ag-zz.github.io/}
\\
\sphinxbottomrule
\end{tabulary}
\sphinxtableafterendhook\par
\sphinxattableend\end{savenotes}

\sphinxAtStartPar
Im Rahmen der DHd\sphinxhyphen{}Versammlung stellten sich einige Arbeitsgruppen vor und gaben dabei einen konkreten Einblick in ihre Arbeitsweise und ihren Output: von der Buchherausgabe der AG Spiele über Vorträge und Forschungsbeiträge der AG Greening DH bis hin zu Paneldiskussionen, Workshops und einer Summer School bei den AGs Theorie und Zeitungen \& Zeitschriften. Ergänzt wurde dies durch die AG Corner in den Postersessions, die noch einmal kompakter die Bandbreite der AG\sphinxhyphen{}Aktivitäten sichtbar machte. Insgesamt hat mir dieser Überblick gezeigt, wie viel wichtige und vielfältige Arbeit in den AGs geleistet wird – und ich nehme mir vor, mich künftig stärker einzubringen.







\renewcommand{\indexname}{Index}
\printindex
\end{document}